\documentclass[11pt]{article}

\usepackage[margin=1in]{geometry}
\usepackage{amsmath,amssymb,amsthm}
\usepackage{mathtools}
\usepackage{listings}
\usepackage{xcolor}
\usepackage{enumitem}
\usepackage{tikz}
\usetikzlibrary{arrows.meta,positioning,calc}
\usepackage{hyperref}
\hypersetup{colorlinks=true,linkcolor=black,urlcolor=blue,citecolor=black}

\newtheorem{theorem}{Theorem}
\newtheorem{lemma}{Lemma}
\newtheorem{proposition}{Proposition}
\newtheorem{conjecture}{Conjecture}
\theoremstyle{definition}
\newtheorem{definition}{Definition}
\newtheorem{remark}{Remark}
\newtheorem{example}{Example}

\newcommand{\enc}[1]{[\![ #1 ]\!]}
\newcommand{\wbisim}{\approx}
\newcommand{\Nil}{\mathbf{0}}
\newcommand{\writhe}{\operatorname{w}}
\newcommand{\cost}{\operatorname{c}}
\newcommand{\Gate}{\mathsf{Gate}}
\newcommand{\charf}[1]{\chi_{#1}}
\newcommand{\sep}{\mathbin{\vert}}
\newcommand{\ssep}[1]{\mathbin{\vert_{#1}}}
\newcommand{\Hq}{\mathrm{H}}
\newcommand{\fn}{\operatorname{fn}}

\lstdefinelanguage{rholang}{
  morekeywords={new,in,for,contract,def,match,if,else,Nil,true,false},
  sensitive=true,
  morecomment=[l]{//},
  morestring=[b]",
}
\lstset{
  language=rholang,
  basicstyle=\ttfamily\small,
  keywordstyle=\color{blue!70!black}\bfseries,
  commentstyle=\color{green!45!black}\itshape,
  stringstyle=\color{red!60!black},
  columns=fullflexible,
  keepspaces=true,
  showstringspaces=false,
  frame=none,
  xleftmargin=1.5em,
}

\title{Classifying Knot Diagrams with a Spatial--Behavioural Logic\\
\large A logic obtained for free from the encoding of knots as processes}
\author{L. G. Meredith\\ F1R3FLY.io}
\date{\today}

\begin{document}
\maketitle

\begin{abstract}
A companion note encoded knot diagrams as cost-accounted rholang processes and
showed the encoding invariant under the Reidemeister moves up to weak
bisimulation. Weak bisimulation, however, is nearly blind on a closed knot: the
encoding of a closed diagram is a live $\tau$-network with no free names, so its
observable behaviour carries almost no information. This note takes up the other
half. Because the host is a graph-structured lambda theory, it comes with a
logic --- generated from its own constructors and rewrite rules by the OSLF
family of algorithms --- and we get that logic for free the moment we encode
knots into it. Nothing here is designed for knots. We assemble the pieces the
template supplies, add two features it does not yet have (a greatest fixed point
and a graded separating conjunction), and observe that in a reflective calculus
the name predicates are considerably stronger than in a calculus of atomic names,
because a name is a quoted process and a predicate may therefore see into its
structure. We then work a stable of examples: split diagrams, nugatory crossings
and reducedness, connected sums, Conway spheres, alternation, and the writhe.
Two of these are worth the price of admission on their own. Reducedness is a
one-line separation statement, because a nugatory crossing is exactly a crossing
whose deletion disconnects the term --- a property no behavioural logic can
express. And the conjunction \emph{reduced} $\wedge$ \emph{alternating} carves
out precisely the fragment on which the cost ledger of the companion note stops
being a framing and becomes a knot invariant. Throughout we are candid that these
are properties of \emph{diagrams}; we argue this is the right level, since it is
the level at which classical knot theory states the Tait conjectures.
\end{abstract}

\tableofcontents

\section{Introduction}

\subsection{What we are claiming, and what we are not}

The companion note \cite{knots} encodes a knot diagram as a process. A crossing
becomes a gate: the over-strand passes a signal through and produces a private
latch; the under-strand joins on its incoming signal and that latch; chirality is
which strand plays which role. Driving the construction from a
Dowker--Thistlethwaite code makes every arc a single-producer/single-consumer
channel, so a diagram is a flat parallel composition of gates plus injected
current. Taken as a monoidal functor from tangles to processes modulo weak
bisimulation, the encoding is invariant under the Reidemeister moves.

That result is one-directional. It says non-isotopic diagrams are never
\emph{wrongly} separated. It says nothing about whether isotopy classes are
\emph{distinguished}, and there the news is bad: closing a tangle into a knot
restricts every boundary name, so a closed encoding is a live, periodic
$\tau$-network. Weak bisimulation abstracts exactly the internal steps. As an
observable object a closed knot is nearly featureless, and one should not expect
its weak-bisimulation class to carry knot type.

The claim of this note is that the discriminating content is not gone; it has
moved. It is in the \emph{spatial} structure of the term --- how many gates there
are, which arcs they share, how the network falls apart when a gate is removed,
what the chirality pattern is as one walks the traversal. All of that is invisible
to a modal logic and completely visible to a spatial one. The encoding did not
destroy the information; weak bisimulation is simply the wrong instrument for
reading it.

What we are \emph{not} claiming is that any of this is a knot invariant on the
nose. Almost every formula below is a property of a diagram, and we take
Section~\ref{sec:diagrams} to say plainly which is which. We regard this as the
honest and also the useful level: it is the level at which Tait stated his
conjectures, and the level at which the classical results about alternating
diagrams live.

\subsection{The logic is not designed; it is generated}

The reason the logic costs nothing is worth stating up front, because it is the
whole architectural point and is easy to mistake for a slogan.

A graph-structured lambda theory (GSLT) is a triple $(\Sigma, E, R)$: a signature
which is in general a lambda theory, a set of equations imposing a structural
congruence, and an adjoined graph of rewrite rules \cite{omnibus}. Given such a
presentation, the OSLF family of algorithms generates a logic from it: one
structural connective per term former of $\Sigma$, one modality per rewrite rule
of $R$ together with a choice of redex position, and a propositional layer
stratified by the categorical strength available. The generation does not know
which theory it is applied to.

The consequence for us is immediate. The companion note encodes knots into the
rho calculus. The rho calculus is a GSLT with a published presentation. So the
logic for knot diagrams is whatever the generator emits on that presentation ---
we do not get to choose it, and equally we do not have to build it. What the
present note contributes is not a logic but the observation that the generated
one already says the things a knot theorist wants said, plus two additions to the
template (\S\ref{sec:additions}) and one observation about reflection
(\S\ref{sec:names}) that make the saying practical.

\begin{remark}[Why the base theory is the encoding]
One could instead present the tangle category $\mathbf{Tang}$ directly as a GSLT
and generate a logic from \emph{its} constructors and moves. The companion note
declines that route and so do we, for the reason given there: the discriminating
material is the gate --- which strand produces the latch, whether the latch is
matched or counted, how gates sit beside one another under $\sep$ --- and gates
are objects on the process side. A direct presentation of the tangle category has
no gates in it to talk about. It also keeps the resolution reading of
\cite[\S9]{knots} in play, since the four-port hole is what a gate becomes when
its latch is forgotten (\S\ref{sec:relax}).
\end{remark}

\subsection{Ideal logic, checkable fragment}\label{sec:ideal}

One methodological commitment governs everything below. We distinguish the
\emph{ideal} logic --- adequate for the idealised calculus, and the thing that
\emph{defines} the invariant --- from \emph{fragments}, restricted so that model
checking terminates, which \emph{decide} it. These are different artifacts with
different jobs, and it is a mistake to ask one to do the other's work.

The precedent is familiar and should reassure the topologist. The Kauffman
bracket is \emph{defined} by a sum over $2^n$ resolutions; nobody computes it that
way, and nobody thinks the definition is impugned by the existence of a better
algorithm. Likewise here: the adequate logic says what the invariant is, and a
user chasing a terminating check restricts the theory, or the connectives, or
both, and accepts that the restricted logic no longer resolves everything. Since
the spatial and behavioural connectives are generated from the theory, restricting
the theory restricts them automatically; the connectives arising from the
collection semantics --- infinitary joins, and hence the fixed point of
\S\ref{sec:nu} --- can force non-termination independently and may need
restricting on their own account. A logic whose model checking terminates is
necessarily not adequate for the whole of a Turing-complete host.

We flag which side of this line each example falls on as we go.

\section{The logic, assembled}\label{sec:logic}

\subsection{What the template supplies}

We recall only what we use. The rho calculus has term formers
$\Nil$, $*n$ (drop), $n!(Q)$ (output), $\mathtt{for}(x \leftarrow n)P$ (input),
parallel composition, and $@P$ (quote); parallel composition carries an
associative--commutative equational theory, and the communication rewrite is the
single interaction rule. Feeding that presentation to the generator yields three
layers.

\paragraph{Structural connectives.}
One type former $f^\sharp$ per term former $f$, whose inhabitants are the terms
whose root former is $f$. We write $\mathsf{out}(n)\varphi$ for the predicate
holding of $n!(Q)$ with $Q \models \varphi$, and
$\mathsf{in}(n)\varphi$ for the analogous predicate on receipts. The interesting
one is parallel composition. Because $\sep$ is associative--commutative, the
generated connective is a genuine \emph{separating conjunction} in the sense of
spatial and separation logics \cite{caires-cardelli,reynolds}:
\[
P \models \varphi \sep \psi
\quad\Longleftrightarrow\quad
P \equiv Q \sep R \ \text{ with } \ Q \models \varphi,\ R \models \psi .
\]
Everything in \S\ref{sec:examples} that a modal logic cannot say is said with
this connective.

\paragraph{Behavioural modalities.}
One primitive modality $\langle K_j\rangle^{\vec x :: \vec A}B$ per rewrite rule
and choice of redex position, where $K_j$ is the one-hole context determined by
that position. These are \emph{context-labelled}: the label is the context in
which the step fires, not merely an action name. We return to why that matters in
\S\ref{sec:caires}.

\paragraph{Propositional layer.}
Conjunction and $\top$ come from the finite-limits fragment; disjunction,
implication and negation require more, and a specification supplies them or does
not. We assume conjunction, disjunction, and negation throughout the ideal logic,
and say so when a fragment gives one up.

\paragraph{Restriction.}
The encoding of a closed diagram is $\mathtt{new}\ a_1,\dots,a_{2n}\ \mathtt{in}\
\{\cdots\}$, so every arc name is bound. We need a connective that gets under the
binder. Following the spatial-logic tradition we take a hidden-name quantifier:
\[
P \models \Hq x.\,\varphi
\quad\Longleftrightarrow\quad
P \equiv \mathtt{new}\ x\ \mathtt{in}\ Q \ \text{ with } \ Q \models \varphi,
\ x \ \text{fresh}.
\]

\begin{remark}[Restriction hides names from observation, not from description]
\label{rem:crux}
This is the crux of the whole note and deserves to be isolated. Restriction makes
a name unobservable: no context can communicate on it, which is exactly why a
closed knot's behaviour is all $\tau$ and why weak bisimulation sees nothing. It
does not make the name \emph{undescribable}. $\Hq x.\varphi$ reaches under the
binder and speaks about the structure so hidden. The spatial layer therefore
recovers precisely what closing the diagram took away from the behavioural layer.
The companion note's pessimism about closed encodings is a fact about
$\wbisim$, not about the term.
\end{remark}

\subsection{Names are processes}\label{sec:names}

In a calculus of atomic names, a name predicate can do two things: test equality
and test freshness. Caires' spatial logic for the $\pi$-calculus \cite{caires} is
of that kind, and its name apparatus --- name equality, the revelation
connective, the fresh-name quantifier --- is exactly what atomic names permit.

The rho calculus is reflective. Names are not atoms; a name is a quoted process,
with $@$ and $*$ mediating between the two sorts and $@(*n) = n$ as the only
equation. A name predicate may therefore \emph{see into the structure of the
name}:
\[
n \models @\varphi
\quad\Longleftrightarrow\quad
n = @P \ \text{ with } \ P \models \varphi .
\]
This is namespace logic \cite{namespace}: a full logic on names, in which one
describes a \emph{namespace} by a property rather than enumerating it. The
difference from the atomic setting is not cosmetic, and it pays three times here.

\paragraph{Arc names carry their traversal index.}
The parametric builder of \cite[App.~A]{knots} constructs arc $k$ of a knot as the
reflected name $@[*\mathit{knot}, k]$: a structured name whose components are the
diagram's unforgeable token and the arc's position in the DT traversal. A name
predicate can therefore recover both. ``$a$ is an arc of \emph{this} diagram'' and
``$a$ is the arc at traversal position $k$'' are name-level statements, not
data-level ones. In an atomic-name setting the index would have to be carried as a
message and the logic could speak of it only by observing traffic --- which, for a
closed knot, is precisely what is unavailable.

\paragraph{Walking the traversal is a name operation.}
Consequently the walk that the alternation formula of \S\ref{sec:alt} performs ---
from arc $k$ to arc $k+1$ --- can be taken at the name level rather than by
searching the term for the gate that consumes $a_k$. The same property has a
spatial reading and a nominal reading; the nominal one is dramatically cheaper to
check, since it replaces a search over $n$ gates by a computation on a name.

\paragraph{A name may quote a tangle.}
Most consequentially for \S\ref{sec:mutation}, a name can quote an arbitrary
process, hence an arbitrary sub-tangle. ``The channel $c$ carries a tangle
satisfying $\varphi$'' is then simply $c \models @\varphi$ --- a statement about
the name, with no communication required. This is the same mechanism by which
Rholang~1.4 types a channel by the language definition it carries
\cite[\S9]{omnibus}, and it is the natural way to talk about what sits on the far
side of a Conway sphere.

\begin{remark}[Two boundaries, one mechanism]
It is worth noticing that reflection is doing for names what the separating
conjunction does for processes: both are ways of describing a decomposition
without executing it. $@\varphi$ opens a name to reveal the process inside;
$\varphi \sep \psi$ opens a process to reveal the parts. A knot diagram is a
structure made entirely of such decompositions, which is why the two connectives
between them cover so much of the classical vocabulary.
\end{remark}

\subsection{Two additions to the template}\label{sec:additions}

The generated logic as it stands lacks two things we need. Both are independent
features rather than knot-specific hacks, and both should probably migrate into
the rubric of \cite{omnibus} rather than living here; we develop them here so that
this note is self-contained.

\subsubsection{A greatest fixed point}\label{sec:nu}

Add to the grammar a propositional variable $X$ and the former $\nu X.\varphi$,
with $X$ occurring only positively in $\varphi$. The semantics is the standard
one and requires nothing exotic: a formula denotes a collection of witnesses, the
collections are ordered by inclusion, $\varphi$ positive in $X$ induces a monotone
operator on that order, and $\nu X.\varphi$ denotes its greatest post-fixed point,
which exists by Knaster--Tarski provided the order has arbitrary joins.

Two comments on the cost. First, in the vocabulary of \cite{splitting} the demand
for arbitrary joins is a constraint on the \emph{collection} component, not on
the structural or behavioural ones --- the fixed point is paid for by the shape of
the container in which witnesses are aggregated. That is also why $\nu$ is the
first casualty of a terminating fragment: infinitary collection connectives break
termination independently of how thin one makes the theory fragment.

Second, and specific to us, the addition is not optional. The gate of \cite{knots}
is a recursive definition that reinstates itself at the latch join, so a gate's
own behaviour is infinite and no finite modal formula pins it. The characteristic
formula of \S\ref{sec:chi} is necessarily a $\nu$-formula. This is a satisfying
place for the feature to enter, since it enters on the strength of an ordinary
design decision about the object language rather than as apparatus imported to
make a particular example work; alternating diagrams (\S\ref{sec:alt}) then come
along as a second use of something already bought.

\begin{remark}[What $\nu$ buys is uniformity]\label{rem:uniform}
It is worth being precise about the gain, because it is easy to overstate. For a
\emph{fixed} diagram with $2n$ arcs, alternation is a finite conjunction over the
arcs and needs no fixed point at all. What $\nu$ provides is a \emph{single}
formula, the same for every diagram, in place of a $2n$-fold conjunction that must
be rewritten for each $n$. Uniformity, not raw expressive power over a fixed
input. This matters for the same reason it matters in program logics: a
classifier one has to regenerate per input is not a classifier.
\end{remark}

\subsubsection{Graded separation}\label{sec:graded}

The separating conjunction splits a term. For knots we care about \emph{how much}
of the interface the two halves share, since that number is precisely the
classical tangle-decomposition hierarchy. Define, for a set $\vec a$ of restricted
names,
\[
P \models \varphi \ssep{k} \psi
\quad\Longleftrightarrow\quad
P \equiv \Hq \vec a.\,(Q \sep R)
\ \text{ with } \
Q \models \varphi,\ R \models \psi,\ \bigl|\fn(Q) \cap \fn(R) \cap \vec a\bigr| = k .
\]
$\ssep{k}$ is derived rather than primitive: it is $\sep$ under $\Hq$, with the
cardinality condition expressed by the name predicates of \S\ref{sec:names}, which
can say of a name that it is an arc of this diagram and can therefore count the
shared ones. We treat it as a defined abbreviation throughout.

\section{Characteristic formulae}\label{sec:chi}

\subsection{The idea, imported}\label{sec:caires}

Caires \cite{caires} identifies the notion we need: a \emph{characteristic
formula} for a process --- a single formula satisfied by exactly the processes
equivalent to it. Characteristic formulae turn an equivalence question into a
model-checking question, which is the move a classifier is made of.

We import the idea and not the logic, and the distinction matters. Caires'
modalities are labelled by actions. Ours are labelled by \emph{contexts}
(\S\ref{sec:logic}), which is a strictly finer instrument: the label records the
one-hole context in which a step fires, so two steps that perform the same action
in different surroundings are distinguished. For gates this is exactly what one
wants, because chirality --- the whole content of a crossing --- is a fact about
which half of the gate sits in which context. In an action-labelled logic one
recovers chirality by encoding it into a name discipline and then reasoning about
names; here it is read directly off the label. Add to this the stronger name
predicates of \S\ref{sec:names}, and the two logics differ in both of the
ingredients a characteristic formula is built from.

\subsection{The gate}

Recall the gate:
\begin{lstlisting}
def Gate(oi, oo, ui, uo) = {
  new l in {
      for (s <- oi)          { oo!(*s) | l!(Nil) }
    | for (s <- ui & _ <- l) { uo!(*s) | Gate(oi,oo,ui,uo) }
  }
}
\end{lstlisting}
Its characteristic formula describes a restriction of a two-way separation, with
the recursive occurrence caught by the fixed point:
\[
\charf{\Gate}(oi,oo,ui,uo)
\;\coloneqq\;
\nu X.\ \Hq \ell.\Bigl(
  \mathsf{in}(oi)\bigl[\mathsf{out}(oo)\top \sep \mathsf{out}(\ell)\top\bigr]
  \ \sep\
  \mathsf{in}(ui,\ell)\bigl[\mathsf{out}(uo)\top \sep X\bigr]
\Bigr).
\]
Read left to right: the term is, up to structural congruence, the restriction of a
name $\ell$ over a parallel composition of exactly two receipts; the first is a
receipt on $oi$ whose continuation separates into an output on $oo$ and an output
on $\ell$; the second is a join on $ui$ and $\ell$ whose continuation separates
into an output on $uo$ and something again satisfying $X$. The greatest fixed
point is the right one: we want the gate to be reinstated forever, not to be
well-founded.

Two features of the encoding are legible in the formula and are worth naming.
Chirality is the asymmetry between the two conjuncts --- which one produces
$\ell$ and which joins on it --- so $\charf{\Gate}$ and the formula for the mirror
crossing differ by swapping $oi \leftrightarrow ui$, $oo \leftrightarrow uo$. And
the \emph{matched} latch of \cite{knots} is visible as the placement of $X$: the
gate reinstates inside the under-branch, so each latch belongs to one
instantiation. Had the two receipts been made persistent instead --- a
\emph{counting} latch --- $X$ would not appear and the formula would be smaller
and weaker.

\subsection{The relaxation lattice}\label{sec:relax}

$\charf{\Gate}$ sits at the top of a lattice of progressively weaker formulae,
obtained by forgetting parts of it. This is where a classifier's dials live, and
the lattice has a pleasing bottom.

\begin{center}
\begin{tabular}{lll}
\hline
formula & forgets & what remains \\
\hline
$\charf{\Gate}$ & nothing & the gate up to $\wbisim$ \\
$\charf{\Gate}^{-\nu}$ & the reinstatement & one passage through the crossing \\
$\charf{\Gate}^{\pm}$ & the continuations & chirality only \\
$\charf{\Gate}^{\circ}$ & the latch $\ell$ & a four-port hole \\
\hline
\end{tabular}
\end{center}

The bottom row is the point. Forget the latch entirely and what is left is a
predicate on four ports with no over/under content --- which is exactly the
resolution-context of \cite[\S9]{knots}, the hole into which the two Kauffman
smoothings are plugged. The dynamic gate and the Kauffman smoothing are therefore
not material from two unrelated halves of a subject; they are the top and bottom
of one relaxation lattice, and the bracket is what one computes after relaxing all
the way down. We do not develop the bracket here, but we note that the axis exists
and that the classifier and the state sum are graded by the same parameter.

\section{A stable of examples}\label{sec:examples}

Throughout, $D$ is a diagram with $n$ crossings, $\enc{D}$ its encoding, and
$a_1,\dots,a_{2n}$ its arcs, with $o(i)$ and $u(i)$ the passes at which crossing
$i$ is met as over- and under-strand.

\subsection{Graded separation}\label{sec:sep}

One scheme covers four classical notions. Because the encoding is a flat parallel
composition of gates, and because the arcs are the only shared names, a
decomposition of the term is a decomposition of the diagram, and the number of
shared arcs is the number of points in which the separating curve meets it.

\begin{center}
\begin{tabular}{cll}
\hline
shared arcs & spatial statement & classical notion \\
\hline
$0$ & $\enc{D} \models \varphi \ssep{0} \psi$ & split link \\
$0$, after deleting one gate & $\enc{D} \models \charf{\Gate}(\cdots) \sep (\varphi \ssep{0} \psi)$ & nugatory crossing \\
$2$ & $\enc{D} \models \varphi \ssep{2} \psi$ & connected sum \\
$4$ & $\enc{D} \models \varphi \ssep{4} \psi$ & Conway sphere; site of mutation \\
\hline
\end{tabular}
\end{center}

The correspondence is not an analogy. A simple closed curve in the plane meeting
the diagram transversally in $k$ points cuts the arc set into two parts sharing
exactly $k$ arcs, and conversely a $k$-shared decomposition of the term exhibits
such a curve. The classical hierarchy --- split, composite, mutable --- is a
hierarchy in $k$, and the logic is graded by $k$ because the separating
conjunction is graded by its interface.

\begin{remark}[This is what a modal logic cannot say]
Bisimulation cannot count crossings, and it cannot see that a network falls into
two halves that do not communicate. Two diagrams whose encodings are weakly
bisimilar may have different $k$-decompositions for every $k$. Separation is
where the spatial layer earns its name, and it is a fortunate accident --- or
rather, not an accident at all --- that knot theory is largely a subject about
separating curves.
\end{remark}

\subsection{Reduced diagrams}\label{sec:reduced}

The cleanest instance, and the one to check by hand. A crossing is
\emph{nugatory} if there is a simple closed curve meeting the diagram
transversally at that crossing and nowhere else; a diagram is \emph{reduced} if it
has no nugatory crossing. Reducedness is the standing hypothesis of most of the
classical theorems about diagrams, and it is usually explained with a picture.

In the encoding it is a separation statement. Delete gate $i$ from $\enc{D}$. The
traversal cycle breaks into two paths,
\[
\pi^{+}_i = [\,a_{o(i)},\ \dots,\ a_{u(i)-1}\,],
\qquad
\pi^{-}_i = [\,a_{u(i)},\ \dots,\ a_{o(i)-1}\,],
\]
the second being the latch gap of \cite[\S5]{knots}. Crossing $i$ is nugatory
exactly when no remaining gate has ports in both paths --- when, that is, the rest
of the term separates with no shared arcs.

\begin{proposition}[Reducedness is a separation formula]\label{prop:reduced}
Crossing $i$ of $D$ is nugatory if and only if
\[
\enc{D} \ \models\
\Hq \vec a.\Bigl(\charf{\Gate}(a_{o(i)-1},a_{o(i)},a_{u(i)-1},a_{u(i)})
\ \sep\ \bigl(\varphi \ssep{0} \psi\bigr)\Bigr)
\]
for some $\varphi, \psi$, where $\varphi$ and $\psi$ have their free arcs in
$\pi^{+}_i$ and $\pi^{-}_i$ respectively. Hence $D$ is reduced iff the negation of
this holds for every $i$, a formula of size linear in $n$.
\end{proposition}

\begin{proof}[Proof sketch]
Left to right: a curve meeting $D$ only at crossing $i$ separates the other $n-1$
crossings into those inside and those outside, and since the curve crosses no arc,
no remaining gate has ports on both sides; take $\varphi, \psi$ to be the
$\charf{\Gate}$-conjunctions of the two families. Right to left: a $0$-shared
decomposition of $\enc{D} \setminus G_i$ exhibits two families of gates with
disjoint arc support, whose union is all arcs but those at crossing $i$; the two
paths $\pi^{\pm}_i$ are then separated by a curve through $i$ alone.
\end{proof}

We stress how little machinery this used. Non-separability is the primitive the
connective was built on; reducedness is that primitive applied to the term with
one gate removed. Nothing was encoded, and no picture was needed.

\subsection{Alternating diagrams}\label{sec:alt}

$D$ is \emph{alternating} if, following the traversal, one passes alternately over
and under. In the encoding, arc $a_k$ is produced by pass $k$ and consumed by pass
$k+1$, so alternation says: whenever $a$ is the over-out port of a gate it is the
under-in port of a gate, and conversely.

Write $\mathsf{ov}(a) \coloneqq \exists b,c,d.\ \charf{\Gate}(a,b,c,d) \sep \top$
for ``$a$ is the over-in port of some gate'', and $\mathsf{un}(a)$ for the
analogous under-in statement. The walk alternates, and is a greatest fixed point
over the pair:
\[
\begin{aligned}
A_{O}(a) &\;=\; \exists b,c,d.\ \bigl(\charf{\Gate}(a,b,c,d) \sep \top\bigr) \wedge A_{U}(b),\\
A_{U}(a) &\;=\; \exists b,c,d.\ \bigl(\charf{\Gate}(c,d,a,b) \sep \top\bigr) \wedge A_{O}(b),
\end{aligned}
\qquad
\mathsf{Alt} \;\coloneqq\; \nu (A_O, A_U).\ A_O(a_1) \vee A_U(a_1).
\]
In the first clause $a$ is the over-in port, so the pass consuming $a$ is an
over-pass and its out-arc is $b$; the next pass must then be an under-pass, which
is $A_U(b)$. The second clause is the mirror.

\begin{remark}[The recursion is bounded by the diagram]\label{rem:bounded}
The traversal closes after $2n$ steps, so the fixed point is exact at its $2n$-th
unfolding: $\nu^{2n} = \nu$ here, not an approximation. Alternation is therefore
\emph{finitely} checkable per diagram in spite of $\nu$ being an infinitary
connective in general. This is the good case, and it is worth understanding why it
is good: the recursion follows the \emph{structure} --- the arc cycle --- rather
than the reduction. Recursions that follow the reduction are the unbounded ones.
By \S\ref{sec:names} the walk can be taken on names, so the check is $O(n)$ name
operations rather than a search over gates.
\end{remark}

\subsection{Tait, and where the ledger becomes an invariant}\label{sec:tait}

Now the two previous formulae are conjoined, and the payoff is not a curiosity.

$\mathsf{Red} \wedge \mathsf{Alt}$ is exactly the hypothesis of the Tait
conjectures. By the theorem of Kauffman, Murasugi and Thistlethwaite
\cite{kauffman,murasugi,thistlethwaite}, a reduced alternating diagram realises
the crossing number of its link. So a diagram satisfying the conjunction carries a
certificate of minimality: a formula in the generated logic, checkable in time
linear in $n$ by Remark~\ref{rem:bounded} and Proposition~\ref{prop:reduced},
implies a statement about the knot and not merely the diagram.

The sharper payoff concerns the companion note. There the cost ledger is shown to
satisfy $\cost(\enc{D}) = 2\writhe(D)$ identically, and the Perko pair is used to
insist that this is a \emph{framed} invariant: two reduced $10$-crossing diagrams
of one knot with different writhes and therefore different ledgers. Little's 1900
claim that the writhe of a reduced diagram is a knot invariant is false, and the
Perko pair is the counterexample.

But Little's claim is true on the alternating fragment. By the flyping theorem of
Menasco and Thistlethwaite \cite{flyping}, any two reduced alternating diagrams of
the same link are related by flypes, and a flype preserves the writhe. Hence:

\begin{proposition}[The formula carves out the fragment]\label{prop:fragment}
On diagrams satisfying $\mathsf{Red} \wedge \mathsf{Alt}$, the cost ledger
$\cost(\enc{D}) = 2\writhe(D)$ is an invariant of the knot type, not merely of the
framing.
\end{proposition}

This is the demonstration we would lead with. A spatial--behavioural formula ---
one conjunct from the separating conjunction, one from the fixed point, neither
designed for the purpose --- identifies exactly the region on which a quantity
from the \emph{other} note upgrades from a framed invariant to an unframed one.
The Perko pair, which is non-alternating, is outside the region, exactly as it
must be. Two independently developed pieces of apparatus agree about where the
boundary lies, and neither was tuned to make them agree.

\subsection{The writhe as a counting formula}\label{sec:writhe}

The writhe itself can be written as a formula, but only after one is careful about
where the sign lives.

The sign $\varepsilon_i$ of a crossing is \emph{not} its chirality: chirality is
which strand goes over, whereas the sign needs that together with both strand
directions \cite[Rem.~2]{knots}. So the plain gate term does not contain the sign,
and no formula over $\enc{D}$ can count the writhe. In the cost-accounted theory
$C(\mathrm{rho})$, however, the token stack is a sort of the grammar
\cite[\S6]{omnibus} and each crossing's located stack carries $\varepsilon_i$.
The sign is then part of the term, hence visible to the structural layer, and
\[
\mathsf{W}_{w} \;\coloneqq\;
\bigvee_{\substack{p+m = n \\ p - m = w}}
\underbrace{\charf{\Gate}^{+} \sep \cdots \sep \charf{\Gate}^{+}}_{p}
\ \sep\
\underbrace{\charf{\Gate}^{-} \sep \cdots \sep \charf{\Gate}^{-}}_{m}
\ \sep\ \top
\]
says that the diagram has writhe $w$, where $\charf{\Gate}^{\pm}$ is the gate
formula conjoined with the predicate on its located stack's sign.

The interest is not the formula, which is routine, but the cross-check. The same
quantity now has two independent readings --- a ledger reading, obtained by
running the process and summing debits, and a logical reading, obtained by
inspecting the term without running it --- and they must agree. That they do is a
small piece of evidence that the encoding and the generated logic are describing
the same object, and cheap demonstrations of two routes to one answer are worth
more to a sceptical reader than an additional theorem.

\subsection{Mutation: the open case}\label{sec:mutation}

The $k=4$ row of \S\ref{sec:sep} is the interesting risk, and we record it as an
open problem rather than a result.

Conway mutation cuts a diagram along a sphere meeting it in four points, rotates
the tangle inside, and reglues. It is the classical source of hard-to-separate
pairs: the Kinoshita--Terasaka and Conway knots are mutants, share a Jones
polynomial, share an Alexander polynomial, and are distinguished only by
considerably heavier machinery.

In our setting a mutation is a re-plugging of a four-arc subterm --- and by
\S\ref{sec:sep} the site is \emph{visible}: $\enc{D} \models \varphi \ssep{4}
\psi$ locates it, and by \S\ref{sec:names} the tangle on the far side can be named
by a quotation and described by $@\varphi$ without communicating with it. Whether
the induced equivalence \emph{separates} mutants is a different question, and it
is open.

\begin{conjecture}[Mutants]\label{conj:mutants}
There is a formula in the graded-separation fragment satisfied by the
Kinoshita--Terasaka diagram and not by the Conway diagram.
\end{conjecture}

We flag this as the natural first computation for the classifier, precisely
because it can come out either way and one would want to know which before
building further on it. If it holds, it is a place where the process invariant
strictly beats the polynomial, on the same graded-separation scheme as the easy
cases of \S\ref{sec:reduced}; if it fails, the failure locates a definite limit
and is worth as much.

\section{Diagrams, not knots}\label{sec:diagrams}

Every formula above is a predicate on diagrams. $\mathsf{Alt}$ detects an
alternating \emph{diagram}; ``$K$ is an alternating knot'' is an existential over
the diagrams of $K$. $\mathsf{Red}$, the $\ssep{k}$ statements, and
$\mathsf{W}_w$ are all likewise diagram-level. We have said this once per example
and say it once more in general, because the companion note had to correct exactly
this confusion in another guise --- reading the writhe, a diagram quantity, as a
knot quantity --- and we would rather not reintroduce it in new clothes.

We regard the level as correct rather than as a limitation, for three reasons.

First, it is where the classical theory lives. Reducedness, alternation, flypes,
the crossing number of a \emph{diagram}, the Tait conjectures --- these are all
diagram-level notions, and the theorems that make them valuable are theorems about
diagrams that have knot-level consequences. Section~\ref{sec:tait} is precisely of
that shape.

Second, the passage from diagram to knot is a quotient, and quotients are
downstream work. A diagram predicate $\varphi$ induces two knot predicates,
$\exists D \in K.\, \varphi(D)$ and $\forall D \in K.\, \varphi(D)$, and which one
is wanted depends on the application. The interesting mathematics is in showing
that a particular $\varphi$ has one of these forms computable --- which for
alternation is a genuine theorem and not a formality.

Third, the honest statement of what the companion note proved is also
diagram-level in this sense. Reidemeister invariance says the encoding does not
distinguish diagrams it should not. The classifier's job is the converse, and the
converse is naturally attacked by finding diagram properties fine enough to
separate, then asking which of them descend.

\section{Checkability}\label{sec:check}

Following \S\ref{sec:ideal}, we record what a terminating fragment would keep and
give up.

\paragraph{Keep.}
The structural connectives, including $\sep$ and the derived $\ssep{k}$; the name
predicates of \S\ref{sec:names}; the context-labelled modalities; conjunction and
disjunction; and $\nu$ \emph{under a fixed unfolding bound}. By
Remark~\ref{rem:bounded} the recursions this note actually uses follow the
traversal and are exact at $2n$ unfoldings, so the bound costs nothing on these
examples.

\paragraph{Give up.}
The composition adjunct $\varphi \triangleright \psi$. Validity for spatial logics
with the adjunct is where undecidability bites hardest, and the generated
structural layer of \cite{omnibus} does not currently supply a right adjoint to
$\sep$ in any case. We therefore place the adjunct in the ideal logic only and
exclude it from the checkable fragment explicitly, rather than leaving its status
implicit. Negation is the other candidate: it is what adequacy needs and what
decidability least likes, and $\mathsf{Red}$ as stated in
Proposition~\ref{prop:reduced} uses it. Since the negation there is applied to a
formula of bounded size over a fixed finite diagram, it is harmless; we note that
this is a property of the use and not a general licence.

\paragraph{What terminates.}
On a fixed diagram with $n$ crossings, $\mathsf{Alt}$ is $O(n)$ name operations,
$\mathsf{Red}$ is $n$ separation tests each over a term of $n$ gates, and
$\mathsf{W}_w$ is a single pass over the located stacks. The $\ssep{4}$ search of
\S\ref{sec:mutation} is the expensive one, since the number of four-element arc
interfaces is $\binom{2n}{4}$; whether it admits a better algorithm is open and is
the practical question standing between this note and a usable tool.

\section{Open problems}

\begin{enumerate}[leftmargin=1.4em]
\item \textbf{Mutants.} Conjecture~\ref{conj:mutants}: does the
graded-separation fragment separate the Kinoshita--Terasaka and Conway diagrams?
This is the first computation to do.
\item \textbf{Descent.} Which of the diagram predicates of
\S\ref{sec:examples} have computable existential or universal closures over an
isotopy class? Alternation is the model case and is a theorem; reducedness is the
next one to ask about.
\item \textbf{Faithfulness.} The companion note reformulates the converse of
its invariance theorem as the question whether the encoding \emph{reflects} as
well as preserves context-labelled transition structure. The present note's
formulae are candidate separating witnesses; assembling them into a proof, or
into a counterexample, is the main open problem behind both notes.
\item \textbf{The relaxation lattice and the bracket.}
Section~\ref{sec:relax} exhibits $\charf{\Gate}^{\circ}$ --- the gate with its
latch forgotten --- as the Kauffman resolution-context. Making the passage from
$\charf{\Gate}$ to $\charf{\Gate}^{\circ}$ into a statement relating the
classifier and the state sum, rather than an observation that both live on one
axis, is the most promising structural question here.
\item \textbf{Migration.} The fixed point of \S\ref{sec:nu} and the graded
separation of \S\ref{sec:graded} are independent of knots and belong in the
rubric of \cite{omnibus}. Whether $\ssep{k}$ should be primitive there or remain
derived is a question about the rubric, not about this application.
\end{enumerate}

\begin{thebibliography}{99}

\bibitem{knots} L.~G.~Meredith. \emph{Knots as cost-accounted rholang processes: a
  Dowker--Thistlethwaite-driven encoding and weak-bisimulation invariance.}
  F1R3FLY.io, 2026.

\bibitem{omnibus} \emph{Graph-structured lambda theories: a reference for the
  working developer.} F1R3FLY.io, 2026.

\bibitem{splitting} L.~G.~Meredith. \emph{Splitting the mind of Zeus:
  realizability, collection monads, and the choice content of the McBride
  derivative in a finite-limits metalanguage.} F1R3FLY.io, 2026.

\bibitem{caires} L.~Caires. \emph{Behavioral and spatial observations in a logic
  for the $\pi$-calculus.} In Foundations of Software Science and Computation
  Structures (FoSSaCS), LNCS 2987, Springer, 2004.

\bibitem{caires-cardelli} L.~Caires and L.~Cardelli. \emph{A spatial logic for
  concurrency (part I).} Information and Computation \textbf{186}(2):194--235,
  2003.

\bibitem{namespace} L.~G.~Meredith and M.~Radestock. \emph{Namespace logic: a
  logic for a reflective higher-order calculus.} In Trustworthy Global Computing,
  LNCS 3705, Springer, 2005.

\bibitem{rho} L.~G.~Meredith and M.~Radestock. \emph{A reflective higher-order
  calculus.} Electronic Notes in Theoretical Computer Science
  \textbf{141}(5):49--67, 2005.

\bibitem{reynolds} J.~C.~Reynolds. \emph{Separation logic: a logic for shared
  mutable data structures.} In Logic in Computer Science (LICS), 2002.

\bibitem{oslf} Operational Semantics in Logical Form (OSLF): a family of
  algorithms auto-generating logics and type systems from graph-structured lambda
  theories.

\bibitem{kauffman} L.~H.~Kauffman. \emph{State models and the Jones polynomial.}
  Topology \textbf{26}:395--407, 1987.

\bibitem{murasugi} K.~Murasugi. \emph{Jones polynomials and classical conjectures
  in knot theory.} Topology \textbf{26}:187--194, 1987.

\bibitem{thistlethwaite} M.~B.~Thistlethwaite. \emph{A spanning tree expansion of
  the Jones polynomial.} Topology \textbf{26}:297--309, 1987.

\bibitem{flyping} W.~Menasco and M.~Thistlethwaite. \emph{The classification of
  alternating links.} Annals of Mathematics \textbf{138}:113--171, 1993.

\bibitem{conway} J.~H.~Conway. \emph{An enumeration of knots and links, and some
  of their algebraic properties.} In Computational Problems in Abstract Algebra,
  Pergamon, 1970.

\bibitem{lickorish} W.~B.~R.~Lickorish. \emph{An introduction to knot theory.}
  Graduate Texts in Mathematics 175, Springer, 1997.

\end{thebibliography}

\end{document}
